\documentclass[11pt,a4paper]{article}
\usepackage[utf8]{inputenc}
\usepackage[vietnamese]{babel}
\usepackage{amsmath,amssymb}
\usepackage{geometry}
\geometry{a4paper, total={180mm,260mm}, left=15mm, right=15mm, top=18mm, bottom=18mm}

\usepackage{tikz}
\usetikzlibrary{shapes.geometric, arrows.meta, positioning, calc}
\tikzset{
  arr/.style={-{Stealth[scale=0.8]}},
  every node/.append style={align=center}
}
\usepackage{tcolorbox}
\tcbuselibrary{skins,breakable}
\usepackage{fontawesome}
\usepackage{fancyhdr}
\usepackage{enumitem}

\pagestyle{fancy}
\fancyhf{}
\renewcommand{\headrulewidth}{0.5pt}
\renewcommand{\footrulewidth}{0.5pt}
\fancyhead[L]{\small\bfseries\color{violet!85!black}\faBook\ SỔ TAY KIẾN THỨC TOÁN 2}
\fancyhead[R]{\small\bfseries\color{violet!85!black}Lớp Toán Thầy Lĩnh -- loptoanthaylinh.vn}
\fancyfoot[C]{\small\color{gray}-- Sổ tay kiến thức: Bài 27. Thực hành gấp, cắt, ghép, xếp hình và vẽ đoạn thẳng --}

\setlength{\parindent}{0pt}
\setlength{\parskip}{0pt}

\begin{document}

% --- BANNER BÀI HỌC ---
\begin{tcolorbox}[
  enhanced,
  colback=violet!8!white,
  colframe=violet!70!black,
  arc=3mm, boxrule=1.2pt,
  left=3mm, right=3mm, top=2mm, bottom=2mm,
  halign=center
]
  {\LARGE\bfseries\color{violet!90!black}\faBookmark\ BÀI 27: THỰC HÀNH GẤP, CẮT, GHÉP, XẾP HÌNH VÀ VẼ ĐOẠN THẲNG}\\[1mm]
  {\itshape\footnotesize\color{black!75} Kĩ năng dùng thước vẽ đoạn thẳng đúng kích thước, ghép hình sáng tạo và đếm hình trên 1 trang A4}
\end{tcolorbox}

\vspace{-2mm}

\noindent
\begin{minipage}[t]{0.49\textwidth}\vspace{0pt}
\begin{tcolorbox}[
  enhanced,
  colback=violet!4!white,
  colframe=violet!75!black,
  arc=2mm, boxrule=0.8pt,
  left=2mm, right=2mm, top=1.8mm, bottom=1.8mm,
  title={\footnotesize\bfseries\color{white}\faKey\ 1. Cách vẽ đoạn thẳng có độ dài cho trước}
]
{\footnotesize
\begin{itemize}\setlength{\itemsep}{1pt}\setlength{\leftskip}{-2.5mm}
  \item \textbf{Bài toán:} Vẽ đoạn thẳng $AB$ dài $6\text{ cm}$.
  \item \textbf{Bước 1:} Đặt thước kẻ lên giấy, chấm một điểm $A$ trùng với vạch số \textbf{0} của thước.
  \item \textbf{Bước 2:} Tìm vạch số \textbf{6} trên thước và chấm điểm $B$.
  \item \textbf{Bước 3:} Dùng bút chì nối từ điểm $A$ đến điểm $B$ men theo mép thước.
  \item \textbf{Bước 4:} Nhấc thước ra, ghi tên đoạn thẳng $AB$ và số đo $6\text{ cm}$.
\end{itemize}
}
\end{tcolorbox}
\end{minipage}
\hfill
\begin{minipage}[t]{0.49\textwidth}\vspace{0pt}
\begin{tcolorbox}[
  enhanced,
  colback=teal!4!white,
  colframe=teal!75!black,
  arc=2mm, boxrule=0.8pt,
  left=2mm, right=2mm, top=1.8mm, bottom=1.8mm,
  title={\footnotesize\bfseries\color{white}\faSliders\ 2. Ghép hình và đếm hình}
]
{\footnotesize
\begin{itemize}\setlength{\itemsep}{1pt}\setlength{\leftskip}{-2.5mm}
  \item \textbf{Ghép hình:}\\
  - Ghép $2$ hình tam giác vuông bằng nhau được $1$ hình chữ nhật hoặc $1$ hình tam giác lớn.\\
  - Ghép các que tính thành hình tam giác ($3$ que), hình vuông ($4$ que bằng nhau), hình chữ nhật ($6$ que).
  \item \textbf{Kĩ năng đếm hình thông minh:}\\
  - Đánh số thứ tự từng mảnh ghép nhỏ ($1, 2, 3,\dots$).\\
  - Đếm hình đơn $\to$ Đếm hình ghép từ 2 mảnh $\to$ Đếm hình ghép toàn thể.
\end{itemize}
}
\end{tcolorbox}
\end{minipage}

\vspace{-1.5mm}
\begin{tcolorbox}[
  enhanced,
  colback=orange!4!white,
  colframe=orange!80!black,
  arc=2mm, boxrule=0.8pt,
  left=2mm, right=2mm, top=1.5mm, bottom=1.5mm,
  title={\footnotesize\bfseries\color{white}\faSitemap\ 3. Minh họa thao tác vẽ đoạn thẳng AB dài 5 cm trên thước kẻ}
]
\centering
\begin{tikzpicture}[x=1.3cm, y=0.8cm]
  % Thước kẻ
  \draw[thick, fill=yellow!20] (0,0) rectangle (6,1);
  \foreach \x in {0,...,5} {
    \draw[thick] (\x, 1) -- (\x, 0.6) node[below, font=\tiny\bfseries] {\x};
    \ifnum\x<5
      \draw (\x+0.5, 1) -- (\x+0.5, 0.75);
    \fi
  }
  \node[font=\tiny\bfseries] at (5.7, 0.4) {cm};

  % Đoạn thẳng vẽ bên trên
  \draw[very thick, violet!90!black] (0, 1.4) -- (5, 1.4);
  \filldraw[red!80!black] (0, 1.4) circle (2.5pt) node[above=2pt, font=\scriptsize\bfseries] {$A$ (vạch 0)};
  \filldraw[red!80!black] (5, 1.4) circle (2.5pt) node[above=2pt, font=\scriptsize\bfseries] {$B$ (vạch 5)};
  \node[font=\scriptsize\bfseries, teal!90!black] at (2.5, 1.7) {Đoạn thẳng $AB = 5\text{ cm}$};
\end{tikzpicture}
\end{tcolorbox}

\vspace{-1.5mm}
\noindent
\begin{minipage}[t]{0.49\textwidth}\vspace{0pt}
\begin{tcolorbox}[
  enhanced,
  colback=green!4!white,
  colframe=green!60!black,
  arc=2mm, boxrule=0.8pt,
  left=2mm, right=2mm, top=1.8mm, bottom=1.8mm,
  title={\footnotesize\bfseries\color{white}\faCheckCircle\ 4. Ví dụ thực hành chuẩn mực}
]
{\footnotesize
\textbf{Đếm hình tam giác:}\\
Cho hình chữ nhật được chia bởi $1$ đường chéo:\\
- Hình đơn: $2$ hình tam giác nhỏ.\\
- Tổng cộng: Có $\mathbf{2}$ hình tam giác.\\
Cho hình vuông có $2$ đường chéo cắt nhau:\\
- Có $4$ hình tam giác đơn.\\
- Có $4$ hình tam giác đôi (ghép từ 2 tam giác đơn).\\
$\implies$ Có tất cả $\mathbf{8}$ hình tam giác!
}
\end{tcolorbox}
\end{minipage}
\hfill
\begin{minipage}[t]{0.49\textwidth}\vspace{0pt}
\begin{tcolorbox}[
  enhanced,
  colback=red!4!white,
  colframe=red!75!black,
  arc=2mm, boxrule=0.8pt,
  left=2mm, right=2mm, top=1.8mm, bottom=1.8mm,
  title={\footnotesize\bfseries\color{white}\faExclamationTriangle\ 5. Mẹo tránh bẫy trắc nghiệm}
]
{\footnotesize
\begin{itemize}\setlength{\itemsep}{1pt}\setlength{\leftskip}{-2.5mm}
  \item \textbf{Bẫy đặt mép thước:}\\
  Nhiều bạn đặt đầu đoạn thẳng trùng với \textbf{mép ngoài cùng của thước} thay vì trùng với \textbf{vạch số 0}. Điều này làm đoạn thẳng bị vẽ hụt độ dài.
  \item Khi đếm hình, luôn ghi lại các cặp hình ghép để không bị sót hoặc đếm trùng lặp.
\end{itemize}
}
\end{tcolorbox}
\end{minipage}


\end{document}
